\chapter{Motivation and Introduction}

\section{Background on Credit Card Fraud Detection}
Credit card fraud is a significant and escalating challenge in the global financial ecosystem. As digital transactions, e-commerce platforms, and contactless payments continue to grow rapidly, malicious actors continuously develop sophisticated methodologies to bypass traditional security measures. Fraudulent typologies have evolved from simple physical card theft to complex vectors including phishing, skimming, account takeover, and card-not-present fraud. 

Synthetic reviews of machine-learning pipelines for payment-card fraud summarise a recurring pattern of operations: heterogeneous data sources (transaction ledgers, device fingerprints, behavioural sequences) are distilled into tabular or sequence features, imbalance-aware learners are tuned under severe computational and latency budgets, and deployable systems face continual adaptation pressure because adversaries explicitly probe decision boundaries rather than behaving as stationary draws from a fixed distribution \cite{lucas_jurgovsky_ccfraud_survey2020,yousefi_fraud_survey2019}. These surveys complement practitioner-oriented handbook material by foregrounding methodological choices—feature taxonomies, model families, imbalance treatments—that recur across empirical studies despite differences in jurisdictions and datasets \cite{fraud_handbook}.

As detailed in fraud analytics literature, building robust defenses requires confronting the algebraic structure of class priors alongside the semantics of transactional variables \cite{fraud_handbook}. Learning under extreme imbalance is not merely a bookkeeping inconvenience: optimisation objectives aligned with naive accuracy incentivise collapsing toward the majority class unless resampling schemes, asymmetric loss formulations, kernel modifications, or other algorithm-level corrections explicitly re-weight the rarity of positives \cite{he_garcia_imbalanced2009}. In parallel, non-stationarity implies that hypotheses validated on retrospective windows degrade when fraud rings shift modality; longitudinal monitoring thus belongs to the same problem statement as predictive accuracy, rather than representing an optional postscript \cite{lucas_jurgovsky_ccfraud_survey2020}.

Consequently, financial institutions rely on adaptive systems that score streams in milliseconds while remaining amenable to audit: the engineering stack must admit traceable scores, reproducible preprocessing, and—when regulators or customers contest an outcome—a defensible reconstruction of why a particular observation fell on the escalate side of policy.

\section{The Machine Learning Black Box Dilemma}
To tackle these complex data challenges, financial institutions overwhelmingly adopt statistical learning algorithms complemented by probabilistic thresholds and monitoring layers. Traditional rule-based systems, which relied on rigid, manually defined thresholds, are increasingly being replaced by dynamic, data-driven approaches. Models such as Gradient Boosting Machines and Deep Neural Networks have become the industry standard due to their ability to process massive volumes of transactional data and identify complex, non-linear patterns.

The tension between fidelity to nonlinear structure and scrutability surfaces in methodological debates beyond fraud: some authors argue that in sufficiently high-impact domains designers should prioritise inherently interpretable model classes when credible, reserving black-box optima for settings where transparency can be relaxed \cite{rudin_stopexplaining2019}. Payment fraud rarely satisfies the latter clause—chargeback exposure, reputational harm, and supervisory expectations push institutions toward models that either remain structurally legible (shallow trees, sparse linear risk scores) or are wrapped in documentation layers that externalise decision logic. In practice, competitive performance often still favours opaque ensembles; the dissertation therefore treats post-hoc explanation as a controlled epistemic interface between those ensembles and human investigators, without equating interface outputs with causal discovery.

However, these highly complex models inherently act as opaque ``black boxes'' in the sense that internal composition maps are not directly human-auditable. They may output calibrated or merely rank-inducing scores; either way, the map from raw features to decisions is not legible without auxiliary analysis. From a business and regulatory perspective, predictive accuracy alone is insufficient. Automated decision support is subject to scrutiny when it materially affects data subjects: the GDPR, for example, foregrounds rights to meaningful information about logic involved in decisions that significantly impact individuals, alongside safeguards when solely automated processing produces legal or similarly significant effects \cite{gdpr_text}. If a legitimate transaction is falsely declined, the institution must be able to justify the decision in terms analysts and, where applicable, customers can interrogate. Operating systems that lack traceable rationale introduces business risk, erodes trust, and complicates compliance audits even when offline metrics appear strong.

\section{Explainable AI and Post-Hoc Analysis}
To resolve the tension between high predictive performance and transparency requirements, this thesis engages Explainable AI with emphasis on \emph{post-hoc} methods as \emph{families} whose evaluation requires explicit benchmarks about human utility, fidelity, or robustness---not surrogate visual appeal alone---as retrospective position papers codify \cite{doshi_velez_kim_interpretability_2017}. Survey-level taxonomies nevertheless distinguish problems (e.g., understanding global behaviour versus explaining a single decision), black-box types (parametric, nonparametric, mixed), and desired explanation forms (rules, attributions, counterfactuals); matching method to problem is itself a design choice rather than a universal recipe \cite{guidotti_blackbox_survey2018,xai_survey}. Post-hoc techniques treat the trained predictor as a queryable oracle: explanations are constructed from input--output behaviour (and sometimes from gradients or structural probes), not necessarily from a human-readable parameterisation of the hypothesis class \cite{guidotti_blackbox_survey2018,molnarInterpretableMachineLearningBook}.

Representative formulations include Local Interpretable Model-agnostic Explanations (LIME), which fits locally weighted interpretable surrogates, and SHapley Additive exPlanations (SHAP), which redistribute prediction deviations using cooperative-game constructions \cite{lime_original,shap_original}. Their appeal is modularity across model classes, yet perturbation-based explainers can be manipulated when evaluated on neighbourhoods that stray from realistic transaction support \cite{slack2020_fooling_lime_shap}. Accordingly, Chapter~5 labels the ExplainableFraud dashboard's heuristic influence mask distinctly from calibrated LIME/SHAP artefacts. The dissertation does \emph{not} treat surrogate vectors as causal attributions independent of preprocessing contracts; summaries remain contractually constrained by neighbourhood semantics, estimator variance, and asymmetric costs downstream.

\section{Rare Events, Ranking Geometry, and RQ1 Preview}
Precision--recall summaries and ROC-based summaries answer different questions about classifier behaviour when class priors skew heavily toward negatives \cite{davis_goadrich_icml06,fraud_handbook}. Roughly speaking, ROC geometry emphasises discrimination between labelled positives and negatives at arbitrary score thresholds; precision--recall geometry makes the scarcity of positives explicit in the denominator of precision and therefore reacts more visibly when escalation policies chase rare fraud at scale. Formal analysis shows that ROC and precision--recall spaces are invertibly related via the confusion matrix yet induce different notions of dominance and visual skew when negatives overwhelm positives \cite{davis_goadrich_icml06}. RQ1 (Chapter~3) revisits threshold selection against this geometric background: scalar accuracy is an inadequate surrogate when the contingency table is overwhelmingly populated by true negatives attainable via trivial majors.

\section{Scientific standpoint and delineation}
This dissertation adopts a \emph{methodological-realist} orientation: stochastic classifiers deployed in supervised settings carve out lawful regularities contingent on labelled archives, preprocessing commitments, and stationarity horizons, yet explanatory overlays remain inferential artefacts whose authority is bounded by explicit contracts \cite{rudin_stopexplaining2019,xai_survey}. The contribution is accordingly triangulated rather than monocausal—bridging probabilistic forecasting, asymmetric evaluation geometry, and interpretability epistemology while refusing to imply that modest reference software alone closes operational risk gaps prevalent in multinational card networks.

The boundary conditions deserve emphasis. Sociotechnical critiques of Explainable AI stress that dashboards can camouflage asymmetric power distributions between model owners and scrutinised subjects; fraud analytics partially inverts this tension because flagged merchants and cardholders seldom access internal investigator consoles \cite{xai_survey,gdpr_text}. Even so, European data-protection instruments continue to incentivise substantive transparency narratives when purely automated profiling affects natural persons---a motive for pairing numeric scores with traceable methodological prose even when PCA masks obscure semantics \cite{gdpr_text}.

Methodologically adjacent areas—deep sequence models, graph-augmented representations, reinforcement policies for sequential decision thresholds—sit outside the analytic core to preserve argumentative sharpness around tabular stochastic classification and attribution masks. Subsequent researchers may remap the scaffolding onto richer feature geometries without overturning foundational claims about imbalance, calibration, or post-hoc limitations.

\section{Research Questions}
To anchor the analytical narrative, three research questions (RQs) are investigated:
\begin{description}
    \item[RQ1] Which evaluation and decision-theoretic arguments justify using precision--recall and threshold-based operating points when fraud prevalence is orders of magnitude smaller than legitimacy?
    \item[RQ2] How should post-hoc explainability theories (local surrogate approximation and additive feature attributions grounded in cooperative games) inform what fraud analysts may legitimately infer from model outputs?
    \item[RQ3] Under which structural assumptions can supervised tabular learners produce probability estimates compatible with qualitative explanation interfaces, despite inevitable distributional drift in adversarial domains?
\end{description}
The methodological substance of each RQ is developed in Chapter~2 (domain theory) and Chapter~3 (proposed mathematical framing). Chapter~4 walks through the ExplainableFraud reference application; Chapter~5 treats empirical evaluation templates and numerical reporting. Software and metrics are illustrative rather than exhaustive production certification.

\section{Thesis Objectives and Contributions}
The dissertation prioritises conceptual clarity and formal reasoning about imbalance \cite{he_garcia_imbalanced2009,fraud_handbook}, evaluation geometry appropriate to rare positives \cite{davis_goadrich_icml06}, probabilistic thresholds, and the epistemic status of post-hoc interpretability claims \cite{guidotti_blackbox_survey2018,gdpr_text}. Bridging predictive strength with explanatory accountability engages both cooperative-game attributions and local surrogates \cite{lime_original,shap_original,xai_survey}.

The main contributions emphasise theory-aware methodology:
\begin{enumerate}
    \item A structured exposition of imbalance-aware evaluation semantics and operational threshold selection (\emph{versus} naive accuracy optimisation).
    \item A unified discussion of complementary post-hoc explanation paradigms: local surrogate fidelity (LIME), additive Shapley-style attributions (SHAP), and their interplay with constrained financial monitoring contexts.
    \item A methodological architecture that couples predictive scoring $f_\theta$, explicit decision thresholds, and attributable evidence vectors grounded in reproducible preprocessing contracts.
    \item A proof-of-concept implementation (Chapter~4) illustrating how algebraic feature vectors align with reproducible artefacts, HTTP/DTO boundaries, and analyst-facing explanation views.
\end{enumerate}

\section{Document Structure}
The remainder of this dissertation unfolds as follows. Chapter~2 reviews state-of-the-art fraud analytics and Explainable AI, emphasising taxonomy, imbalance metrics, limitative properties of explanations, and regulatory motivation. Chapter~3 formalises the proposed methodological solution: stochastic classification model, discriminative surrogates, attribution layers, and operational constraints independent of incidental implementation choices. Chapter~4 presents the ExplainableFraud companion implementation: user workflow, architecture, API contracts, dependency injection, and artefacts such as screenshots and reference listings. Chapter~5 concentrates on reproducible experimentation, metric reporting templates, and critique of applicability. Chapter~6 summarises conclusions, limitations, and research outlook. Supplemental appendices consolidate notation, conceptual glossaries, threshold sketches, and explanation quality heuristics for readers rotating between algebraic, empirical, and software-facing passages.

Reading strategies may vary by examiner specialty without violating argumentative dependencies: governance reviewers might traverse Chapters~2--3 before auditing Chapter~4 lineage hooks; empiricists may prioritise Chapter~5 prior to revisiting Chapter~3 formalisms whenever notation demands clarification; practitioners evaluating reproducibility can alternate Chapter~4 (\emph{how}) with appendix symbolism (\emph{vocabulary}). Anchoring vignettes appear sparingly—preference attaches to transferable methodological scaffolding rather than institution-specific folklore unavailable within anonymised corpora.